\section{Related work}\label{sec:related}

\paragraph{Viability kernels and robust invariance.}
The viability kernel is Aubin's \cite{Aubin1991}; its discrete-time robust form, the largest set from which a target tube can be held against a bounded adversary, is the minimax reachability problem of Bertsekas and Rhodes \cite{BertsekasRhodes1971,Bertsekas1972}, surveyed with the invariance literature by Blanchini \cite{Blanchini1999}. Saint-Pierre's algorithm \cite{SaintPierre1994} computes kernels by backward set recursion on a grid, which is the computation we specialise to one dimension, where the sets are intervals and the recursion has closed forms. The exact bounded-defect kernel of Theorem~\ref{thm:exact} is the robust kernel of the map perturbed by an additive bounded disturbance; its relation to the Lipschitz inflation, Proposition~\ref{prop:lip}, is the standard inner-approximation argument, with equality for affine maps wherever the backward recursion binds. Regnier and De Lara \cite{RegnierDeLara2015} develop robust viability for a harvested ecosystem under set-valued uncertainty and are the closest antecedent for the floor-class formulation; De Lara and Doyen's monograph \cite{DeLaraDoyen2008} and the stochastic viability of Doyen and De Lara \cite{DoyenDeLara2010} and Rougé et al.\ \cite{Rouge2013} give the probabilistic counterpart we report alongside the kernels.

\paragraph{Viability in fisheries and groundwater.}
Viability analysis has been applied to bio-economic fishery models by Béné et al.\ \cite{Bene2001}, to recovery paths by Martinet et al.\ \cite{Martinet2007}, to ecosystem-based management by Doyen et al.\ \cite{Doyen2012} and to by-catch by Krawczyk et al.\ \cite{Krawczyk2013,KrawczykPharo2013}; Oubraham and Zaccour \cite{Oubraham2018} survey the field. In groundwater, Pereau and co-authors \cite{Pereau2018,Pereau2019} contrast viability with optimality in aquifer management with environmental flows, and the safe-yield debate \cite{Alley1999,Alley2004,Sophocleous1997,Sophocleous2000,Devlin2005} is the institutional background to the question of which pumping level holds a head. These studies compute kernels or viable paths for a given model; the present paper adds the certification layer (declared defect, finite horizon, retention rule) and the scoring protocol, and tests both on assessment records rather than on calibrated illustrative models.

\paragraph{Management procedures and triggers.}
The management-procedure approach in fisheries \cite{Butterworth2007,Punt2016} evaluates candidate rules by closed-loop simulation against operating models; our closed-loop replays are its deterministic single-model instance, and our kernels are what simulation cannot supply---a guarantee over a class. The precautionary framework of Fisheries and Oceans Canada \cite{DFO2009} defines the limit reference point we certify, from the 2016 assessment \cite{DFO2016}; the 2024 reassessment \cite{Regular2025} redefines it, and Section~\ref{sec:cod-ident} reports the frame under both. The collapse of the stock is documented by Hutchings and Myers \cite{Hutchings1994}, Myers and Cadigan \cite{Myers1995} and Myers et al.\ \cite{Myers1997}; the catch reconstruction is Schijns et al.'s \cite{Schijns2021}. For the Edwards Aquifer the trigger architecture is set out by the Edwards Aquifer Authority \cite{EAA2024}, the Habitat Conservation Plan \cite{EARIP2021} and its National Research Council review \cite{NRC2015}; the recharge series is the U.S.\ Geological Survey's \cite{Umphres2025,Puente1978}, the discharge accounting the Authority's \cite{EAA2023} and the J-17 record the Texas Water Development Board's \cite{TWDB}. Drought triggers as stochastic decision rules are analysed by Steinemann \cite{Steinemann2003}, robust water-resource solutions by Watkins and McKinney \cite{Watkins1997}, and the limits of equivalent-porous-medium models of the karst by Scanlon et al.\ \cite{Scanlon2003}. The Lucas critique \cite{Lucas1976} is the reason a map fitted under one pumping regime is certified only within a declared class.

\paragraph{Companion studies.}
The two forecast-evaluation studies \cite{Abaee2026codforecast,Abaee2026edwforecast} test whether the surplus-production and one-pool maps improve forecasts over naive baselines, which is a different question from the one asked here; the two intervention studies \cite{Abaee2026codintervention,Abaee2026edwintervention} are the registered protocols whose archived campaigns this paper re-derives, and the obstruction calculus \cite{Abaee2026obstruction} treats the incompletely observed case. This paper supersedes the separate reports of the two protocols by giving them a common frame and the exact bounded-defect kernels, which neither had.
